% Figure: draining tank height versus time under Torricelli's law. The
% water height falls as the square-root-of-height outflow removes volume;
% the separable integral gives a height that decays quadratically in time
% and reaches zero at a finite drain time, unlike an exponential decay.
\begin{figure}[htbp]
\centering
\begin{tikzpicture}
\begin{axis}[
  width=11cm, height=6.2cm,
  xlabel={time $t$ (s)}, ylabel={water height $h(t)$ (m)},
  xmin=0, xmax=120, ymin=0, ymax=1.05,
  domain=0:100, samples=200,
  axis lines=left, grid=both,
  major grid style={engblack!12}, font=\small,
  legend pos=north east, legend style={font=\footnotesize},
]
% Torricelli: h(t) = (sqrt(h0) - k t /2)^2 for t < t_drain, clamped to 0.
% h0 = 1, t_drain = 100 so sqrt(h0)/(k/2) = 100 => k/2 = 0.01, k = 0.02.
% h(t) = (1 - 0.01 t)^2 on [0,100].
\addplot[engblue, very thick, domain=0:100]
  {(1 - 0.01*x)^2};
\addlegendentry{Torricelli $h(t)=(\sqrt{h_0}-\tfrac{k}{2}t)^2$}
% comparison: exponential decay with same initial slope, h = exp(-0.02 t)
\addplot[engred, thick, dashed, domain=0:120]
  {exp(-0.02*x)};
\addlegendentry{exponential (same initial slope)}
% finite drain time marker
\addplot[engamber, only marks, mark=*, mark size=1.8pt]
  coordinates {(100,0)};
\node[font=\scriptsize, engamber, anchor=south east]
  at (axis cs:99,0.05) {$t_{\mathrm{drain}}$};
\end{axis}
\end{tikzpicture}
\caption[Draining-tank height under Torricelli's law]{The height of water
in a draining tank under Torricelli's law, where the outflow speed grows
as the square root of the remaining head. Separating the variables and
integrating gives a height $h(t) = (\sqrt{h_0} - \tfrac{k}{2}t)^2$ that
falls to zero at a finite drain time $t_{\mathrm{drain}} =
2\sqrt{h_0}/k$ (amber point), where the tank is empty. An exponential
decay (red dashed) sharing the same initial slope never reaches zero: the
square-root outflow law, weaker than a linear one at low head, still
empties the tank in finite time because the height decays quadratically,
not exponentially. The distinction is the working content of separating
and integrating a nonlinear rate law rather than assuming the
first-order exponential every decay reflexively suggests.}
\label{fig:vol02ch06:tank-drainage}
\end{figure}
